\begin{abstract}
This paper describes StaleyNET, a small system that generates music from a
short text description. We built it from first principles, mainly as a way of
understanding what each part of such a system contributes. All components,
including the GPU kernels, are written in Rust, and the models are trained
from random initialisation on two NVIDIA Jetson devices with about 23 hours of
Creative Commons music. Sound is represented by quantised log-mel
spectrograms, and a decoder-only Transformer learns to predict these tokens
from a caption and, optionally, from lyrics. Because every stage is our own,
we could examine the stages one at a time. These measurements suggest that the
quality of the generated audio is limited mainly by the tokenizer rather than
by the language model: a model trained on a single excerpt reproduces it
exactly, while the tokenizer on its own loses much of the spectral detail.
Guided by this observation, we revised the tokenizer and roughly halved its
reconstruction error. We also describe a simple paired test for whether a
model makes use of the lyrics it is given; our first lyrics-conditioned models
do not yet pass it. Each algorithm is presented with a small worked example,
in the hope that the paper is useful to readers who would like to build a
similar system.
\end{abstract}

\section{Introduction}
\label{sec:introduction}

Generating music from a sentence of text has made remarkable progress in
recent years \cite{agostinelli2023musiclm,copet2023musicgen}. Most current
systems combine several large components that are trained separately: a
neural codec that turns audio into tokens, a text encoder that turns the
prompt into vectors, and a Transformer that relates the two. This division of
labour works very well, although it can make it difficult to tell which
component is responsible for a particular quality, or a particular flaw, in
what one hears.

In this work we approached the problem from the other direction and built a
complete, if modest, system ourselves (\cref{fig:overview}). StaleyNET
computes its spectrograms, learns its tokenizer, trains its Transformer and
reconstructs its waveforms without any pretrained component. It is written in
Rust, and its GPU kernels are compiled from Rust by cuda-oxide \cite{cudaoxide}.
We trained it on an NVIDIA Jetson Orin NX with 16\,GB of memory and on an
NVIDIA Jetson AGX Thor with 128\,GB, two embedded computers designed for
robotics. Owning every stage gave us the opportunity to study each one in
isolation, and much of the paper is about what that study revealed.

\begin{widefigure}
\resizebox{\linewidth}{!}{% Figure "overview", shared by the paper and (as SVG) the book.
\begin{tikzpicture}[figure, node distance=6mm and 7mm]
  \node[block, fill=blockgray, minimum width=22mm] (audio) {recording\\[-1pt]\scriptsize 24\,kHz};
  \node[block, fill=blockblue, minimum width=27mm, right=of audio] (mel) {log-mel spectrogram\\[-1pt]\scriptsize STFT as a matrix product};
  \node[block, fill=blockblue, minimum width=27mm, right=of mel] (rvq) {residual $k$-means\\[-1pt]\scriptsize 4 levels $\times$ 1\,024};
  \node[block, fill=blockblue, minimum width=22mm, right=of rvq] (delay) {delay pattern\\[-1pt]\scriptsize staircase of tokens};
  \node[block, fill=blockviolet, minimum width=27mm, minimum height=13mm, right=of delay] (model) {\textbf{Transformer}\\[-1pt]\scriptsize 6 blocks};
  \node[block, fill=blockamber, minimum width=22mm, anchor=south] (caption) at ($(model.north)+(-13mm,8mm)$) {caption bytes};
  \node[block, fill=blockrose, minimum width=22mm, anchor=south] (lyrics) at ($(model.north)+(13mm,8mm)$) {lyrics bytes};
  \node[block, fill=blockgreen, minimum width=27mm, below=11mm of model] (sample) {sampled tokens\\[-1pt]\scriptsize top-$k$ and guidance};
  \node[block, fill=blockgreen, minimum width=27mm, left=of sample] (decode) {sum of prototypes\\[-1pt]\scriptsize log-mel spectrogram};
  \node[block, fill=blockgreen, minimum width=27mm, left=of decode] (gl) {Griffin-Lim\\[-1pt]\scriptsize phase estimation};
  \node[block, fill=blockgray, minimum width=22mm, left=of gl] (out) {music\\[-1pt]\scriptsize 24\,kHz};
  \draw[flow] (audio) -- (mel);
  \draw[flow] (mel) -- (rvq);
  \draw[flow] (rvq) -- (delay);
  \draw[flow] (delay) -- (model);
  \draw[flow] (caption.south) -- ($(model.north)+(-4mm,0)$);
  \draw[flow] (lyrics.south) -- ($(model.north)+(4mm,0)$);
  \draw[flow] (model) -- node[right, caption text] {one step at a time} (sample);
  \draw[flow] (sample) -- (decode);
  \draw[flow] (decode) -- (gl);
  \draw[flow] (gl) -- (out);
  \node[font=\sffamily\footnotesize\bfseries, text=muted, anchor=south west] at ($(audio.north west)+(0,2.5mm)$) {Training};
  \node[font=\sffamily\footnotesize\bfseries, text=muted, anchor=north west] at ($(out.south west)+(0,-2.5mm)$) {Generation};
\end{tikzpicture}
}
\caption{An overview of StaleyNET. During training (top row), a recording is
converted into a staircase of discrete tokens, which the Transformer learns
to predict from the caption and the lyrics. During generation (bottom row),
the Transformer produces tokens one step at a time, and the tokens are turned
back into sound.}
\label{fig:overview}
\end{widefigure}

Our first models learned text-conditioned structure well above chance, yet the
voices they produced sounded blurred. Two experiments helped us locate the
cause. When we passed held-out recordings through the tokenizer and the
vocoder alone, without any model, most of the spectral detail was already
lost. When we trained the model on a single excerpt, it reproduced that
excerpt token for token. Together these results suggested that the tokenizer
was the main limitation, and a revised tokenizer that quantises groups of
frequency bands separately reduced the reconstruction error by about half.

The contributions of this paper are as follows.
\begin{itemize}[leftmargin=1.4em, itemsep=1pt]
\item A complete text-to-music pipeline written in Rust, including its GPU
  kernels, and trained from scratch on embedded GPUs.
\item A description of each algorithm with a worked example and, where an
  algorithm is easy to get wrong, an argument for its correctness.
\item A way of adding a lyrics input to a trained model without discarding
  what it has learned, together with a paired test of whether the input is
  used.
\item A stage-by-stage analysis that separates the errors of the tokenizer,
  the vocoder and the language model, and a tokenizer study guided by it.
\end{itemize}

We write algorithms in the style of Cormen et al.\
\cite{cormen2022introduction}, as in our earlier tutorial on tiled matrix
multiplication \cite{pathak2026tiled}. Each procedure first states its inputs
and its output in words; its lines are numbered, and the text refers to them.
The range $[a, b)$ denotes the integers $a, \dots, b-1$, indices start at 0, and
losses are measured in nats.

\section{Related work}
\label{sec:related}

AudioLM \cite{borsos2023audiolm} showed that audio can be generated by a
language model over discrete tokens. MusicLM \cite{agostinelli2023musiclm}
conditioned this approach on text, and MusicGen \cite{copet2023musicgen}
simplified it to a single Transformer over the codebooks of a neural codec,
interleaved by a delay pattern. We follow MusicGen in using the delay
pattern, top-$k$ sampling \cite{fan2018hierarchical} and classifier-free
guidance \cite{ho2022cfg}, and we replace its pretrained components with
$k$-means quantisation of log-mel frames and with raw text bytes, as studied
for text by ByT5 \cite{xue2022byt5}.

Jukebox \cite{dhariwal2020jukebox} generated singing from lyrics and learned
the correspondence between words and sound from 1.2 million songs. Singing
synthesis systems such as DiffSinger \cite{liu2022diffsinger} instead work from
a musical score with phoneme durations. Our lyrics experiments rely on the
line-level timings of the JamendoLyrics benchmark
\cite{stoller2019lyrics,durand2023lyrics}. Neural codecs
\cite{zeghidour2021soundstream,defossez2023encodec,kumar2023dac} and neural
vocoders \cite{kong2020hifigan,siuzdak2023vocos} achieve far higher fidelity
than $k$-means tokens and the Griffin-Lim algorithm
\cite{griffin1984signal,perraudin2013fast}, at a training cost beyond our
setting; \cref{sec:pipeline} measures what our simpler choices cost.

\section{From sound to spectrogram and back}
\label{sec:representation}

A waveform sampled at 24\,kHz contains 24\,000 numbers per second, and much of
that information describes the phase of each vibration, to which the ear is
not very sensitive. We therefore describe sound by its log-mel spectrogram:
the energy in 128 frequency bands spaced according to human pitch
perception, measured about 94 times per second and compressed
logarithmically.

\subsection{Resampling}

The training recordings arrive at several sample rates, most often 44.1\,kHz,
and we convert them all to 24\,kHz. Lowering the rate requires first removing
all content above the new Nyquist frequency of 12\,kHz; otherwise that content
folds back into the audible range, an effect known as aliasing.

\begin{definition}[Bandlimited resampling]
Given samples $x_m$ at rate $f_\text{in}$, the resampled signal at rate
$f_\text{out}$ is
\[
y_n \;=\; \sum_{m} x_m \, h\!\left(t_n - m\right), \qquad t_n = n\,\frac{f_\text{in}}{f_\text{out}},
\]
where $h$ is a low-pass filter. Each output sample is a weighted sum of the
input samples near the fractional position $t_n$.
\end{definition}

We use a windowed sinc filter, $h(t) = c\,\operatorname{sinc}(ct)\,w(ct/Z)$, in
which $c$ is the cutoff relative to the input rate, $Z = 24$ is the number of
zero crossings kept on each side, and $w$ is a Kaiser window with $\beta = 8.6$
\cite{smith2002resample}. Evaluating $h$ anew for every output sample would be
slow, but fortunately the fractional part of $t_n$ takes only a few values.

\paragraph{A worked example.} For a conversion from 44\,100\,Hz to 24\,000\,Hz
the two rates share the factor 300, so $t_n = 147n/80$. For $n = 1$ we obtain
$t_1 = 1 + 67/80$. Every $t_n$ has a fraction of the form $p/80$ with
$p \in [0, 80)$, so exactly 80 sets of filter weights are needed. These sets
are called the \emph{phases} of a polyphase filter, and they are computed
once.

\begin{procedure}{\proc{Resample}$(x, f_\text{in}, f_\text{out})$}
{the samples $x_0, x_1, \dots$ of a recording, its sample rate $f_\text{in}$,
and the desired rate $f_\text{out}$, both in Hz.}
{the samples $y_0, y_1, \dots$ of the same sound at rate $f_\text{out}$.}
\li $g = \gcd(f_\text{in}, f_\text{out})$ \\
\li $L = f_\text{out} / g$ \cmt{number of phases} \\
\li $M = f_\text{in} / g$ \cmt{input step per output sample, in $L$-ths} \\
\li $c = 0.97 \cdot \min(1,\, f_\text{out}/f_\text{in})$ \cmt{cutoff} \\
\li $R = \lceil Z / c \rceil$ \cmt{filter half-width} \groupend
\li \kw{for} $p \in [0, L)$ \{ \\
\li \> \kw{for} $j \in [-R, R]$ \{ \\
\li \> \> $h_p[j] = c\,\operatorname{sinc}\!\big(c\,(j - p/L)\big)\; w\big(c\,(j - p/L)/Z\big)$ \\
\li \> \} \\
\li \} \groupend
\li \kw{for} each output index $n$ \{ \\
\li \> $m = \lfloor nM / L \rfloor$ \cmt{whole part of $t_n$} \\
\li \> $p = nM \bmod L$ \cmt{fraction of $t_n$} \\
\li \> $y_n = \sum_{j \in [-R, R]} h_p[j] \; x_{m+j}$ \\
\li \}
\end{procedure}

Lines 1--5 determine the geometry of the filter, lines 6--10 compute the $L$
phases once, and line 14 is the single dot product needed for each output
sample. The factor 0.97 in line 4 keeps the transition band of the filter
just below the new Nyquist frequency. A test confirms that the filter does
its job: a 15\,kHz tone is attenuated by more than 80\,dB when 48\,kHz audio is
converted to 24\,kHz.

\subsection{The spectrogram as a matrix product}

The signal is divided into frames of $N = 1024$ samples, taken every $H = 256$
samples, and each frame is multiplied by a Hann window. The strength of each
frequency in each frame is usually computed with the fast Fourier transform.
We computed it instead with a matrix product, so that the spectrogram runs on
the same GPU kernel as the Transformer and benefits from the same tests
(\cref{sec:gpu}).

\begin{definition}[Log-mel spectrogram]
Let row $f$ of the matrix $F$ hold frame $f$ of the signal, and let $W$ have
entries $W_{nk} = w[n]\,e^{-2\pi i kn/N}$ for $k \in [0, N/2]$, where $w$ is the Hann
window. The log-mel spectrogram is
\[
S \;=\; \ln \max\!\big(|FW|^2 B,\; 10^{-5}\big),
\]
where the square is taken entry by entry and the columns of $B$ are 128
triangular mel filters \cite{slaney1998auditory}.
\end{definition}

Column $k$ of $W$ holds a windowed sine and cosine at $k \cdot 23.4$\,Hz, so row
$f$ of $FW$ measures how much of each frequency frame $f$ contains. The real
and imaginary parts of $W$ are two real matrices, and $B$ adds a third
product. This formulation needs about $N/\log_2 N \approx 100$ times the
arithmetic of an FFT. Even so, on the Orin NX it analyses 60\,s of audio in
47\,ms, against 73\,ms for a multithreaded CPU FFT (\cref{tab:bench-other}),
and it agrees with the FFT to within $10^{-2}$ log units.

\subsection{Estimating the phase}

A generated spectrogram tells us how strong each frequency is, but not its
phase. The Griffin-Lim algorithm \cite{griffin1984signal} estimates a phase by
alternating between two requirements: the magnitudes should match the target,
and the spectrogram should be that of an actual waveform.

\begin{procedure}{\proc{Griffin-Lim}$(\,|S|, T\,)$}
{the target magnitudes $|S|$, one row per frame and one column per frequency,
and the number of rounds $T$.}
{a waveform $x$ whose spectrogram approximates $|S|$.}
\li $X = |S| \cdot e^{i\phi}$ with random phases $\phi$ \\
\li $x = \operatorname{ISTFT}(X)$ \\
\li $P_\text{prev} = \operatorname{STFT}(x)$ \groupend
\li \kw{for} $t \in [0, T)$ \{ \\
\li \> $P = \operatorname{STFT}(x)$ \cmt{nearest consistent spectrogram} \\
\li \> $A = P + \alpha\,(P - P_\text{prev})$ \cmt{momentum} \\
\li \> $X = |S| \cdot A / |A|$ \cmt{keep the phase, restore the magnitude} \\
\li \> $P_\text{prev} = P$ \\
\li \> $x = \operatorname{ISTFT}(X)$ \\
\li \} \groupend
\li \kw{return} $x$
\end{procedure}

Here STFT and ISTFT denote the short-time Fourier transform and its inverse.
Line 5 replaces the current estimate by the spectrogram of a real waveform,
and line 7 restores the target magnitudes. Line 6 is the acceleration
proposed by Perraudin et al.\ \cite{perraudin2013fast}, which extrapolates
along the last change with momentum $\alpha = 0.99$; in the first round it has
no effect, because $P = P_\text{prev}$. We use $T = 64$ rounds, after mapping
the 128 mel bands back to 513 frequencies with a regularised least-squares
inverse of $B$.

\section{From spectrogram to tokens}
\label{sec:tokenizer}

A language model chooses each next symbol from a finite vocabulary, whereas a
spectrogram consists of real numbers. The tokenizer bridges this gap. Each
band of the spectrogram is first standardised with its mean and standard
deviation over the training set, so that quiet and loud bands carry similar
weight. A \emph{step} is then formed from consecutive frames; in our first
tokenizer a step joins 2 frames, giving 256 numbers that cover 21.3\,ms.

\subsection{Residual quantisation}

One can think of residual quantisation as describing a colour first by its
family, then by its shade within the family, and then by finer and finer
corrections. Each level names the closest of $K$ prototypes to whatever the
previous levels have left unexplained.

\begin{definition}[Residual vector quantisation]
Let $E_0, \dots, E_{C-1}$ be codebooks, each a list of $K$ prototype vectors.
A step $x$ is encoded as the indices $c_0, \dots, c_{C-1}$ defined by
\begin{gather*}
r_0 = x, \qquad c_\ell = \argmin_{k \in [0, K)} \lVert r_\ell - E_\ell[k] \rVert^2, \\
r_{\ell+1} = r_\ell - E_\ell[c_\ell],
\end{gather*}
and it is decoded as $\hat x = \sum_\ell E_\ell[c_\ell]$. The vector $r_\ell$ is the
\emph{residual}: the part of $x$ that levels $0, \dots, \ell - 1$ have not yet
described.
\end{definition}

\begin{figure}[htbp]
\centering
% Figure "rvq", shared by the paper and (as SVG) the book.
\begin{tikzpicture}[figure, scale=0.85]
  \fill[blockgray, rounded corners=4pt] (-0.7,-0.7) rectangle (4.7,4.7);
  \foreach \x in {0,...,4} \draw[white, line width=0.8pt] (\x,-0.7) -- (\x,4.7);
  \foreach \y in {0,...,4} \draw[white, line width=0.8pt] (-0.7,\y) -- (4.7,\y);
  \foreach \p in {(0,0),(0,4),(4,4)} \fill[accent!45] \p circle (4pt);
  \fill[accent] (4,0) circle (4.5pt);
  \node[below, font=\sffamily\scriptsize] at (4,-0.1) {$E_0[1] = (4,0)$};
  \draw[->, accent, line width=1.4pt] (0,0) -- (3.93,0);
  \draw[->, teal!70!black, line width=1.4pt] (4,0) -- (3.25,0.75);
  \fill[orange!85!black] (3,1) circle (4pt);
  \node[above, font=\sffamily\scriptsize] at (3,1.08) {$x = (3,1)$};
  \fill[teal!70!black] (3.2,0.8) circle (3pt);
  \node[left, font=\sffamily\scriptsize, teal!60!black] at (3.1,0.55) {$\hat x = (3.2, 0.8)$};
\end{tikzpicture}

\caption{Residual quantisation of $x = (3,1)$ in two levels. Level 0 (blue)
chooses the prototype $(4,0)$, and level 1 (green) adds the correction
$(-0.8, 0.8)$.}
\label{fig:rvq}
\end{figure}

\paragraph{A worked example.} In two dimensions, consider the codebooks
\begin{align*}
E_0 &= \{(0,0),\ (4,0),\ (0,4),\ (4,4)\},\\
E_1 &= \{(0.8,0.8),\ (-0.8,0.8),\ (-0.8,-0.8),\ (0.8,-0.8)\},
\end{align*}
and the point $x = (3,1)$ (\cref{fig:rvq}). The squared distances from $x$ to
the prototypes of $E_0$ are $10, 2, 18, 10$, so $c_0 = 1$ and $r_1 = (-1, 1)$.
The closest entry of $E_1$ to $(-1, 1)$ is $(-0.8, 0.8)$, so $c_1 = 1$ and
$r_2 = (-0.2, 0.2)$. The error falls from $\sqrt 2 \approx 1.41$ after one level to
$\sqrt{0.08} \approx 0.28$ after two.

\subsection{Training the codebooks}

\begin{procedure}{\proc{Train-Codebooks}$(X, C, K, I)$}
{a matrix $X$ whose rows are training steps, the number of levels $C$, the
number of prototypes per level $K$, and the number of rounds $I$.}
{the codebooks $E_0, \dots, E_{C-1}$.}
\li $R = X$ \cmt{what remains to be described} \groupend
\li \kw{for} $\ell \in [0, C)$ \{ \\
\li \> $E_\ell = K$ rows of $R$, chosen at random \groupend
\li \> \kw{repeat} $I$ times \{ \\
\li \> \> $a_i = \proc{Nearest}(R_i, E_\ell)$ for each row $i$ \\
\li \> \> $E_\ell[k] = \operatorname{mean}\{R_i : a_i = k\}$ for each $k$ \\
\li \> \> \kw{if} no row has $a_i = k$ \\
\li \> \> \> $E_\ell[k] = $ a random row of $R$ \\
\li \> \} \groupend
\li \> $R_i = R_i - E_\ell[\proc{Nearest}(R_i, E_\ell)]$ for each row $i$ \\
\li \}
\end{procedure}

\proc{Nearest}$(r, E)$ returns the index of the prototype in $E$ closest to
$r$. Lines 5--6 are the two moves of Lloyd's algorithm \cite{lloyd1982least}:
each row is assigned to its nearest prototype, and each prototype then moves
to the mean of its rows. Lines 7--8 restart a prototype to which no row was
assigned, so that none of the $K$ entries goes unused, and line 10 passes the
residuals on to the next level. We use $I = 20$ rounds on 200\,000 steps drawn
from the training set.

\subsection{The nearest prototype as a matrix product}

\begin{lemma}
For a vector $x$ and prototypes $e_0, \dots, e_{K-1}$,
$\argmin_k \lVert x - e_k \rVert^2 = \argmin_k \big(\lVert e_k \rVert^2 - 2\,x \cdot e_k\big)$.
\end{lemma}
\begin{proof}
Expanding gives $\lVert x - e_k \rVert^2 = \lVert x \rVert^2 - 2\,x \cdot e_k + \lVert e_k \rVert^2$,
and the first term does not depend on $k$.
\end{proof}

In the example above, the squared norms $\lVert e_k \rVert^2$ are $0, 16, 16, 32$
and the dot products $x \cdot e_k$ are $0, 12, 4, 16$, so the scores are
$0, -8, 8, 0$, and the smallest again selects $(4, 0)$. On the GPU the dot
products of 8\,192 steps with 1\,024 prototypes form a single matrix product,
after which a small kernel finds the smallest score in each row.

\subsection{Quantising groups of bands}
\label{sec:band-groups}

The analysis of \cref{sec:pipeline} led us to a refinement. Instead of
quantising a whole step at once, we divide its 128 bands into $G$ equal
groups and give each group its own residual codebooks. With one frame per
step and $G = 4$, each group is a vector of 32 numbers, which 1\,024 prototypes
can describe more finely than a whole step. In our experiments the grouping
helped somewhat less than we expected, and the additional codes helped more
(\cref{sec:appendix-tokenizer}). The codes of a
step are ordered level by level: first the coarse codes of all $G$ groups, then
their first corrections, and so on, so that the coarse-to-fine order relied on
by the delay pattern (\cref{sec:model}) is preserved. With $G = 4$ groups and 4
levels each, a step is described by 16 codes. \Cref{tab:tokenizer-study} in
the appendix compares this and other configurations.

\section{The model}
\label{sec:model}

The model (\cref{fig:model}) is a decoder-only Transformer
\cite{vaswani2017attention,radford2019language}. It reads text and audio as a
single sequence and predicts, at each position, what comes next, much as a
text model predicts the next word.

\begin{figure}[htbp]
\centering
% Figure "model", shared by the paper and (as SVG) the book.
% A pre-normalized Transformer block repeated N times. Every block in the stack
% has the same width, and both residual paths run on the right, well clear of
% the blocks, each leaving the main line at a marked branch point and entering
% its sum from the side.
\begin{tikzpicture}[figure]
  % The input sequence: begin marker, caption bytes, lyrics bytes, separator,
  % then audio positions, each holding C stacked codes.
  \foreach \i/\c in {0/blockgray,1/blockamber,2/blockamber,3/blockamber,4/blockrose,5/blockrose,6/blockrose,7/blockgray}
    \fill[\c, rounded corners=1.5pt] (\i*0.42-2.1,0) rectangle +(0.36,0.36);
  \foreach \i in {8,9,10,11}
    \foreach \l in {0,1,2,3}
      \fill[accent!\the\numexpr 25+15*\l\relax, rounded corners=1pt] (\i*0.42-2.1,\l*0.095) rectangle +(0.36,0.075);
  \node[caption text, anchor=north] at (-1.92,-0.06) {B};
  \node[caption text, anchor=north] at (-1.05,-0.06) {caption};
  \node[caption text, anchor=north] at (0.21,-0.06) {lyrics};
  \node[caption text, anchor=north] at (1.02,-0.06) {S};
  \node[caption text, anchor=north] at (2.1,-0.06) {audio, $C$ codes};

  % The main column.
  \tikzset{stage/.style={block, minimum width=52mm, minimum height=7.5mm}}
  \node[stage, fill=blockgray, minimum width=58mm] (embed) at (0,1.35) {code embeddings $+$ position embeddings};
  \node[stage, fill=blockblue, above=10mm of embed] (norm1) {RMSNorm};
  \node[stage, fill=blockviolet, above=4mm of norm1] (attn) {causal self-attention, $h$ heads};
  \node[circle, draw=inkline, fill=white, inner sep=0pt, minimum size=4mm, above=5mm of attn] (add1) {\scriptsize$+$};
  \node[stage, fill=blockblue, above=10mm of add1] (norm2) {RMSNorm};
  \node[stage, fill=blockgreen, above=4mm of norm2] (ffn) {feed-forward $d \to 4d \to d$, GELU};
  \node[circle, draw=inkline, fill=white, inner sep=0pt, minimum size=4mm, above=5mm of ffn] (add2) {\scriptsize$+$};
  \node[stage, fill=blockblue, above=9mm of add2] (normf) {RMSNorm};
  \node[stage, fill=blockamber, minimum width=58mm, above=4mm of normf] (heads) {$C$ output heads, $K = 1\,024$ scores each};

  % Branch points of the residual paths, midway in the gaps below each norm.
  \coordinate (fork1) at ($(embed.north)!0.5!(norm1.south)$);
  \coordinate (fork2) at ($(add1.north)!0.5!(norm2.south)$);
  \coordinate (lane) at ($(attn.east)+(7mm,0)$);

  \draw[flow] (0,0.45) -- (embed);
  \draw[flow] (embed) -- (norm1);
  \draw[flow] (norm1) -- (attn);
  \draw[flow] (attn) -- (add1);
  \draw[flow] (add1) -- (norm2);
  \draw[flow] (norm2) -- (ffn);
  \draw[flow] (ffn) -- (add2);
  \draw[flow] (add2) -- (normf);
  \draw[flow] (normf) -- (heads);

  \draw[flow, rounded corners=4pt] (fork1) -| (lane |- add1) -- (add1.east);
  \draw[flow, rounded corners=4pt] (fork2) -| (lane |- add2) -- (add2.east);
  \fill[inkline] (fork1) circle (1pt);
  \fill[inkline] (fork2) circle (1pt);

  % The repeated block, and how many times it repeats.
  \node[draw=muted!70, dashed, rounded corners=7pt, inner sep=3.5mm, fit=(fork1) (norm1) (attn) (add1) (norm2) (ffn) (add2) (lane)] (stack) {};
  \node[font=\sffamily\small\bfseries, text=muted, anchor=east] at ($(stack.west)+(-1.5mm,0)$) {$N\times$};
\end{tikzpicture}

\caption{The model. Text bytes and audio tokens share one sequence, and each
audio position carries all the codes of one step. A stack of pre-normalised
blocks, each adding its result to its input, ends in one output head per
code.}
\label{fig:model}
\end{figure}

\subsection{One sequence for text and audio}

The vocabulary contains the 256 possible byte values \cite{xue2022byt5}, three
markers (\code{BEGIN}, \code{SEPARATOR} and \code{PADDING}) and $C \cdot K$ audio
symbols, one for each prototype of each codebook. The text occupies a fixed 128
positions, namely \code{BEGIN}, 126 caption bytes and \code{SEPARATOR}, so that
the audio always begins at the same position. An audio position carries all
$C$ codes of a step, and its input vector is the sum of their embeddings and
the embedding of the position.

\subsection{The delay pattern}

At each position the model predicts all $C$ codes of a step. A finer code,
however, corrects a coarser one and is easier to predict once the coarser one
is known.

\begin{definition}[Delay pattern \cite{copet2023musicgen}]
Position $t$ of the delayed sequence holds, in slot $\ell$, the code
$a_{t-\ell}^{\ell}$ of codebook $\ell$ for frame $t - \ell$, and an empty marker
when $t - \ell$ is not a frame. A piece of $F$ frames becomes $F + C - 1$
positions.
\end{definition}

\begin{figure}[htbp]
\centering
% Figure "delay", shared by the paper and (as SVG) the book.
\begin{tikzpicture}[figure, x=8.5mm, y=7.5mm]
  \definecolor{levelzero}{HTML}{EFE6D2}
  \definecolor{levelone}{HTML}{E4B27C}
  \definecolor{leveltwo}{HTML}{C8692F}
  \foreach \l/\name/\c in {0/a/levelzero,1/b/levelone,2/c/leveltwo} {
    \node[caption text, anchor=east] at (-0.55,-\l) {level \l};
    \foreach \f in {0,...,3} {
      \pgfmathtruncatemacro{\t}{\f+\l}
      \node[block, fill=\c, minimum width=7mm, minimum height=5.8mm, inner sep=0pt] at (\t,-\l) {$\name_{\f}$};
    }
  }
  \foreach \t in {0,...,5} \node[caption text] at (\t,0.75) {\t};
  \node[caption text, anchor=east] at (-0.55,0.75) {position};
  \draw[accent, line width=1pt, rounded corners=4pt] (1.55,0.42) rectangle (2.45,-2.42);
  \node[caption text, text=accent, anchor=north] at (2,-2.5) {predicted together};
\end{tikzpicture}

\caption{The delay pattern for three codebooks and four frames. At position 2
the model predicts $a_2$, $b_1$ and $c_0$, and the coarser codes that they refine
are already part of its input.}
\label{fig:delay}
\end{figure}

In \cref{fig:delay}, position 2 predicts $a_2$, $b_1$ and $c_0$ together. By then
$a_1$, which $b_1$ refines, and $a_0$ and $b_0$, which $c_0$ refines, are already
known. The price is $C - 1$ additional positions, whereas giving each codebook
its own positions would multiply the length of the sequence by $C$.

\subsection{The network and its objective}

In each block, attention lets every position gather information from itself
and from earlier positions, and a small feed-forward network then transforms
each position on its own. Both parts normalise their input
\cite{zhang2019rmsnorm,xiong2020layernorm} and add their result to it:
\begin{align*}
h &= x + \operatorname{Attention}(\operatorname{Norm}(x)), \\
y &= h + \operatorname{FeedForward}(\operatorname{Norm}(h)),
\end{align*}
where $x$ is the sequence entering the block and $y$ the sequence leaving it.
The feed-forward network uses the GELU activation \cite{hendrycks2016gelu}. Our
first model has 6 blocks of width 256 with 8 attention heads (7.1 million
parameters); the larger model of \cref{sec:large} has 8 blocks of width 512
(42.9 million parameters).

At position $p$ the model outputs, for each codebook $\ell$, a probability
distribution over the code $y_{p,\ell}$ that appears at position $p + 1$. The
loss is the cross-entropy, averaged first over the set $V_\ell$ of positions
that have a target in codebook $\ell$ and then over the codebooks:
\[
\mathcal{L} = \frac{1}{C} \sum_{\ell=0}^{C-1} \frac{1}{|V_\ell|} \sum_{p \in V_\ell} -\ln P_\theta\big(y_{p,\ell} \mid z_0, \dots, z_p\big),
\]
where $z_p$ is the input at position $p$ and $\theta$ denotes the parameters. A
model that assigns equal probability to all 1\,024 codes has a loss of
$\ln 1024 \approx 6.93$.

\subsection{Gradients by hand}

As far as we know, no automatic differentiation is available for Rust on the
GPU, so each layer implements its own gradient. Attention is the least
obvious case. For one head, let the rows of $Q$, $K$ and $V$ be the queries,
keys and values of the positions, and let $d$ be the width of the head. Then
\[
S = QK^\top/\sqrt d + M, \qquad P = \softmax(S), \qquad O = PV,
\]
where the mask $M$ is $-\infty$ above the diagonal, so that no position attends
to a later one, and the softmax normalises each row.

\begin{proposition}
Let $\partial O$ be the gradient of the loss with respect to $O$, and similarly
for the other matrices. Then $\partial V = P^\top \partial O$,
$\partial P = \partial O\, V^\top$, $\partial S_{ik} = P_{ik}(\partial P_{ik} - \rho_i)$
with $\rho_i = \sum_j \partial P_{ij} P_{ij}$, $\partial Q = \partial S\, K/\sqrt d$ and
$\partial K = \partial S^\top Q/\sqrt d$.
\end{proposition}
\begin{proof}
The first two identities follow from $O = PV$. Within a row,
$\partial P_{ij}/\partial S_{ik} = P_{ij}(\delta_{jk} - P_{ik})$, so
$\partial S_{ik} = \sum_j \partial P_{ij} P_{ij} (\delta_{jk} - P_{ik}) = P_{ik}(\partial P_{ik} - \rho_i)$.
The last two follow from $S = QK^\top/\sqrt d$. Where $M = -\infty$ we have
$P_{ik} = 0$ and therefore $\partial S_{ik} = 0$.
\end{proof}

Apart from $\partial S$, every gradient is a matrix product and runs on the same
kernel as the forward pass. Because hand-written gradients are easy to get
wrong, a test perturbs three random entries of every parameter tensor by
$\pm 10^{-2}$ and checks that the resulting change of the loss agrees with the
computed gradient to within 5\,\%.

\section{Lyrics}
\label{sec:lyrics}

A caption describes how a piece should sound, but not which words should be
sung. To make the words controllable, the model needs a second text input,
and it needs training examples in which it can hear which words are sung
where.

\subsection{A position map that preserves a trained model}

The lyrics occupy $L = 256$ byte positions between the caption and the
separator. Each position $p$ takes its position embedding from row $\pi(p)$ of
a table, and the choice of $\pi$ determines whether a model trained without
lyrics can later be extended (\cref{fig:lyrics}).

\begin{figure}[htbp]
\centering
\resizebox{\linewidth}{!}{% Figure "lyrics", shared by the paper and (as SVG) the book.
\begin{tikzpicture}[figure, x=1mm, y=1mm]
  % The sequence, in the order the model reads it.
  \node[block, fill=blockgray,  minimum width=9mm]  (b) at (0,0)  {B};
  \node[block, fill=blockamber, minimum width=21mm, right=1.6mm of b] (c) {caption};
  \node[block, fill=blockrose,  minimum width=21mm, right=1.6mm of c] (l) {lyrics};
  \node[block, fill=blockgray,  minimum width=9mm,  right=1.6mm of l] (s) {S};
  \node[block, fill=blockblue,  minimum width=19mm, right=1.6mm of s] (a) {audio};
  % The position table, in row order.
  \node[block, fill=blockgray,  minimum width=9mm,  below=13mm of b] (rb) {0};
  \node[block, fill=blockamber, minimum width=21mm, right=1.6mm of rb] (rc) {1 to 126};
  \node[block, fill=blockgray,  minimum width=9mm,  right=1.6mm of rc] (rs) {127};
  \node[block, fill=blockblue,  minimum width=19mm, right=1.6mm of rs] (ra) {128 to 767};
  \node[block, fill=blockrose,  minimum width=21mm, right=1.6mm of ra] (rl) {768 to 1023};
  \draw[flow, draw=muted] (b) -- (rb);
  \draw[flow, draw=muted] (c) -- (rc);
  \draw[flow, draw=muted] (s.south) to[out=-90, in=90] (rs.north);
  \draw[flow, draw=muted] (a.south) to[out=-90, in=90] (ra.north);
  \draw[flow, draw=accent, line width=0.9pt] (l.south) to[out=-90, in=90] (rl.north);
  \node[caption text, anchor=east] at ($(b.west)+(-1.5,0)$) {sequence};
  \node[caption text, anchor=east] at ($(rb.west)+(-1.5,0)$) {table row};
  \node[caption text, text=accent, anchor=north] at ($(rl.south)+(0,-1)$) {rows added for lyrics};
\end{tikzpicture}
}
\caption{The position map $\pi$. Positions outside the lyrics keep the rows
that they used before the lyrics were added, and the lyrics use new rows
appended to the table.}
\label{fig:lyrics}
\end{figure}

\begin{definition}[Position map]
Suppose the caption occupies positions $[1, 1{+}P)$, the lyrics occupy
$[1{+}P, 1{+}P{+}L)$, and the model without lyrics has $T$ rows. Then
\[
\pi(p) =
\begin{cases}
p & \text{if } p < 1 + P, \\
T + p - 1 - P & \text{if } 1 + P \le p < 1 + P + L, \\
p - L & \text{otherwise}.
\end{cases}
\]
\end{definition}

\begin{lemma}
Every position outside the lyrics uses the row that the same part of the
sequence used without lyrics, and the lyrics use only the new rows
$[T, T + L)$.
\end{lemma}
\begin{proof}
Caption positions are unchanged. The separator and the audio are shifted right
by $L$ positions, and $\pi$ subtracts $L$. The lyrics are mapped onto
$T + [0, L)$.
\end{proof}

For example, with $P = 2$, $L = 3$ and $T = 16$, positions $0$ to $8$ use the rows
\[
0,\ 1,\ 2,\ 16,\ 17,\ 18,\ 3,\ 4,\ 5.
\] A trained model can therefore be extended by
appending $L$ new rows to its table. All learned parameters keep their
meaning, and training simply continues.

\subsection{Pairing audio with the words sung in it}

A training example covers a window of a song, and it should carry the lines
sung in that window. JamendoLyrics \cite{durand2023lyrics} gives the start and
end time of every line of 78 songs in English, French, German and Spanish, and
a window receives each line that overlaps it. Takes recorded with our studio
come with their text but without timings, so we assume that their lines are
spread evenly over the take; this is a fair approximation for short takes of
continuous singing. Because the lyrics corpus is small (4.7 hours, against 23
hours of captioned music), each dataset receives a sampling weight: a weight of
$w$ makes each of its songs $w$ times as likely to be drawn.

\subsection{Testing whether the lyrics are used}
\label{sec:lyrics-test}

A model can lower its loss on the lyrics corpus simply by learning how those
particular songs sound, without reading their words. We therefore test the use
of the lyrics directly, by giving the model the same audio with the correct
words, with no words, and with words from another song.

\begin{procedure}{\proc{Lyrics-Test}$(\theta, \mathcal{D}, B)$}
{a trained model with parameters $\theta$, a set $\mathcal{D}$ of held-out songs
with lyrics, and a number of batches $B$.}
{the mean paired difference $\bar\Delta$, its standard error $s$, and whether
the difference is clearly positive.}
\li \kw{for} $b \in [0, B)$ \{ \\
\li \> draw a batch of windows from the songs in $\mathcal{D}$ \groupend
\li \> $\mathcal{L}^\text{right}_b = $ loss of $\theta$ with each window's own lyrics \\
\li \> $\mathcal{L}^\text{none}_b = $ loss of $\theta$ with empty lyrics \\
\li \> $\mathcal{L}^\text{wrong}_b = $ loss of $\theta$ with lines from another song \\
\li \} \groupend
\li $\Delta_b = \mathcal{L}^\text{wrong}_b - \mathcal{L}^\text{right}_b$ for each $b$ \\
\li $\bar\Delta = \frac{1}{B} \sum_b \Delta_b$ \\
\li $s = \sqrt{\operatorname{Var}(\Delta) / B}$ \groupend
\li \kw{return} $\bar\Delta$, $s$, and whether $\bar\Delta > 2s$
\end{procedure}

Lines 3--5 score the same windows, so the differences in line 7 compare like
with like, and most of the variation between windows cancels. A model that
ignores the words gives $\bar\Delta \approx 0$. A model that reads them should
predict the audio better with the right words than with the wrong ones, giving
$\bar\Delta > 0$, measured in nats per code. The wrong lines in line 5 always
begin at a sung line rather than at silence, and comparing lines 4 and 5 helps
distinguish a reaction to the presence of text from a reaction to its
content.

\section{Generation}
\label{sec:generation}

Generation repeats a single step: the model predicts a distribution for each
codebook at the next position, one code is drawn from each, and the drawn codes
are appended to the sequence.

\subsection{The key/value cache}

Position $p$ attends to the keys and values of positions $0, \dots, p$. These do
not change once they have been computed, because no position attends forward,
so a cache can keep them \cite{pope2023inference}.

\begin{procedure}{\proc{Decode-Step}$(z, p, \mathit{cache})$}
{the codes $z$ to place at position $p$, the position $p$ itself, and the keys
and values of positions $0, \dots, p - 1$ for every block.}
{the scores of every codebook for position $p + 1$; the cache is extended by
position $p$.}
\li $h = \operatorname{embed}(z) + \operatorname{position}(\pi(p))$ \groupend
\li \kw{for} each block $j$ \{ \\
\li \> $q, k, v = $ query, key and value projections of $\operatorname{Norm}(h)$ \\
\li \> $\mathit{cache}_j[p] = (k, v)$ \\
\li \> $h = h + \operatorname{Attend}(q,\, \mathit{cache}_j[0 \,..\, p])$ \\
\li \> $h = h + \operatorname{FeedForward}(\operatorname{Norm}(h))$ \\
\li \} \groupend
\li \kw{return} the output-head scores of $\operatorname{Norm}(h)$
\end{procedure}

\begin{theorem}
Feeding a sequence position by position through \proc{Decode-Step} yields, at
every position, the same scores as a single forward pass over the whole
sequence.
\end{theorem}
\begin{proof}
We argue by induction over the blocks. At block 0, the input at position $p$
depends only on the token and the position of $p$. Suppose that the inputs of
block $j$ agree at all positions up to $p$. Then the keys and values computed in
line 3 agree, and when position $p$ is processed the cache holds exactly those
of positions $0, \dots, p$. In the full pass the mask gives zero weight to later
positions, so attention at $p$ combines the same keys and values with the same
weights, and the feed-forward network acts on each position separately. The
outputs of block $j$, which are the inputs of block $j + 1$, therefore agree up
to position $p$.
\end{proof}

A test checks the theorem numerically, to within $10^{-4}$ at every position,
for models with and without lyrics; the second case also checks the position
map $\pi$. Because the cache has a fixed capacity, a step allocates no memory.

\subsection{Sampling and guidance}

For each codebook we keep the $k = 250$ most probable codes and draw among
them \cite{fan2018hierarchical}. Because the text was removed from 10\,\% of the
training examples, the same model can predict both with the text, giving
scores $z^\text{text}$, and without it, giving $z^\text{none}$. Classifier-free
guidance \cite{ho2022cfg} amplifies their difference:
\[
z = z^\text{none} + s \,\big(z^\text{text} - z^\text{none}\big),
\]
where $s$ is the guidance scale. For instance, if two codes score $(2, 1)$ with
the text and $(1, 1)$ without it, the first has probability 0.73; with $s = 3$
the scores become $(4, 1)$, and its probability rises to 0.95. We use $s = 3$,
which doubles the computation per step. Lyrics longer than $L$ bytes are cut
at line breaks into windows that are presented in turn, each over an equal
share of the output, while the most recent audio is kept as context.

\section{GPU kernels}
\label{sec:gpu}

Matrix products account for almost all of the arithmetic: the spectrogram,
the nearest-prototype search, the projections of the Transformer and both
products inside attention. Our tutorial \cite{pathak2026tiled} derives the
shared-memory tiled product in detail; here we describe one further step,
register blocking, and the way we made one gradient repeatable.

\subsection{Register blocking}

In the tiled kernel, a block of $16 \times 16$ threads copies tiles of $A$ and
$B$ into shared memory, which reduces traffic to global memory sixteenfold
\cite{pathak2026tiled,kirk2022programming}. Every multiply-add still reads two
values from shared memory, however, and shared memory is slower than a
thread's registers. Register blocking \cite{volkov2008benchmarking} lets each
thread compute an $r \times r$ group of outputs, so that every value it reads is
used $r$ times (\cref{fig:blocking}).

\begin{figure}[htbp]
\centering
% Figure "blocking", shared by the paper and (as SVG) the book.
\begin{tikzpicture}[figure, x=4.6mm, y=4.6mm]
  \foreach \r in {0,...,7} \foreach \c in {0,...,7} {
    \pgfmathtruncatemacro{\own}{(mod(\r,4)==1) && (mod(\c,4)==2)}
    \ifnum\own=1 \fill[accent!70, rounded corners=1pt] (\c,-\r) rectangle +(0.88,-0.88);
    \else \fill[blockgray, rounded corners=1pt] (\c,-\r) rectangle +(0.88,-0.88); \fi
  }
  \node[caption text, anchor=south] at (4,0.25) {output tile of one block, $8 \times 8$};
  \node[caption text, anchor=west, align=left] at (8.6,-2.2) {thread $(t_x, t_y) = (2, 1)$ owns\\rows $1, 5$ and columns $2, 6$};
  \node[caption text, anchor=west, align=left] at (8.6,-5.4) {neighbouring threads read\\neighbouring words of \textit{tileB}};
\end{tikzpicture}

\caption{Register blocking with $4 \times 4$ threads and $r = 2$. Each thread
computes $r^2$ outputs spread over the tile, in rows $t_y + 4i$ and columns
$t_x + 4j$. Our kernel uses $16 \times 16$ threads, $r = 4$ and a $64 \times 64$
tile.}
\label{fig:blocking}
\end{figure}

\begin{procedure}{\proc{Blocked-Thread}$(t_x, t_y)$}
{the coordinates $(t_x, t_y)$ of a thread within a block of $16 \times 16$
threads; the block computes a $64 \times 64$ tile of $C = AB$, with $r = 4$.}
{the 16 entries of the tile owned by the thread, written to $C$.}
\li $\mathit{sum}[i][j] = 0$ for $i, j \in [0, r)$ \groupend
\li \kw{for} each phase $t$ \{ \\
\li \> copy a $64 \times 16$ tile of $A$ into \textit{tileA} \\
\li \> copy a $16 \times 64$ tile of $B$ into \textit{tileB} \\
\li \> \proc{Barrier}() \groupend
\li \> \kw{for} $u \in [0, 16)$ \{ \\
\li \> \> $\alpha_i = \mathit{tileA}[t_y + 16i][u]$ for $i \in [0, r)$ \\
\li \> \> $\beta_j = \mathit{tileB}[u][t_x + 16j]$ for $j \in [0, r)$ \\
\li \> \> $\mathit{sum}[i][j] = \mathit{sum}[i][j] + \alpha_i \beta_j$ for $i, j \in [0, r)$ \\
\li \> \} \groupend
\li \> \proc{Barrier}() \\
\li \} \groupend
\li store $\mathit{sum}[i][j]$ at row $t_y + 16i$, column $t_x + 16j$ of the tile
\end{procedure}

\proc{Barrier}() waits until every thread of the block has reached it. Lines 7
and 8 read $2r = 8$ values from shared memory, and line 9 performs $r^2 = 16$
multiply-adds, so each multiply-add costs $2/r = 0.5$ reads instead of 2.
Global traffic falls as well: each phase copies $2 \cdot 64 \cdot 16$ values for
$64 \cdot 64 \cdot 16$ multiply-adds, so the product of an $M \times K$ and a
$K \times N$ matrix fetches
\[
\frac{M}{64} \cdot \frac{N}{64} \cdot \frac{K}{16} \cdot 2 \cdot 64 \cdot 16 = \frac{2MNK}{64}
\]
values, a quarter of the $2MNK/16$ needed with $16 \times 16$ tiles. The two
barriers play the same roles as in the tiled kernel, whose correctness is
proved in \cite{pathak2026tiled}. On the Orin NX the blocked kernel is 1.8 to
2.3 times faster than the tiled one (\cref{tab:bench-gemm}).

\subsection{A repeatable embedding gradient}

The gradient of the embedding table adds the gradient $g_p$ of each position $p$
into the row of its token $z_p$. Many positions share a token. With atomic
additions, threads would add into a shared row in an order that varies from
run to run, and because floating-point addition is not associative, the result
would vary as well. We fix the order with a counting sort.

\begin{procedure}{\proc{Embedding-Gradient}$(z, g, V)$}
{the token $z_p$ and the gradient vector $g_p$ of every position $p$, and the
number of rows $V$ of the embedding table.}
{the gradient $G$ of the embedding table, one row per token.}
\li $\mathit{count}[v] = |\{p : z_p = v\}|$ for $v \in [0, V)$ \cmt{on the CPU} \\
\li $\mathit{start}[v] = \sum_{u < v} \mathit{count}[u]$ \\
\li $\mathit{list} = $ positions sorted by token, then by position \groupend
\li \kw{for} each token $v$ and column $c$, in parallel \{ \cmt{on the GPU} \\
\li \> $G[v][c] = \sum g_p[c]$ over $p$ in $\mathit{list}[\mathit{start}[v] \,..\, \mathit{start}[v+1])$, in order \\
\li \}
\end{procedure}

For the tokens $z = (2, 0, 2, 1)$, the counts are $(1, 1, 2)$, the starts are
$(0, 1, 2, 4)$ and the list is $(1, 3, 0, 2)$, so token 2 receives $g_0 + g_2$ in
that order. Each entry is summed by one thread in a fixed order, so the result
is the same on every run, and a whole training run can be repeated bit for
bit on the same device.

\section{Experiments}
\label{sec:experiments}

The hardware is described in \cref{tab:hardware} in the appendix. The main
corpus is the Song Describer Dataset \cite{manco2023songdescriber}, 706
two-minute Creative Commons excerpts with 1\,106 captions written by
listeners. Songs are divided between training and evaluation by artist, so
that no voice appears on both sides, which leaves 667 songs for training and 39
for evaluation. We report spectral errors as mean absolute differences of
natural-log mel energies; one unit corresponds to about 8.7\,dB.

\paragraph{Correctness and speed.} All 38 tests pass on the Orin NX, and the
kernel tests also pass on the Thor. They compare every kernel with a CPU
implementation at deliberately awkward sizes, check every gradient
numerically, verify the cache theorem, and confirm that a small network can
memorise a fixed batch. On the Orin NX, register blocking makes a training
step of the first model 1.8 times faster than tiling and 9.3 times faster than
the direct kernel (\cref{tab:bench-train}).

\paragraph{First training runs.} \Cref{tab:training} summarises the runs with
the first tokenizer. A run of 1\,500 steps reached a held-out loss of 4.93, and
a run of 15\,000 steps on the Thor reached 4.02. When we restricted the data to
the 281 songs whose captions mention a voice, a model fine-tuned from the
longer run reached its best held-out loss of 4.23 after only 250 steps, after
which its held-out loss rose while its training loss continued to fall; a
model trained from scratch on the same songs reached 4.72. The vocal subset
appears to be too small for long training.

\begin{table}[htbp]
\centering
\caption{Best held-out loss with the first tokenizer, in nats per code
(uniform guess: 6.93).}
\label{tab:training}
\small
\setlength{\tabcolsep}{4pt}
\begin{tabular}{llrr}
\toprule
run & data & steps & loss \\
 & & & (nats per code) \\
\midrule
first run, Orin NX & all songs & 1\,500 & 4.93 \\
longer run, Thor & all songs & 15\,000 & 4.02 \\
fine-tuned, Thor & vocal songs & 250 & 4.23 \\
from scratch, Orin NX & vocal songs & 5\,500 & 4.72 \\
\bottomrule
\end{tabular}
\end{table}

\subsection{What limits the sound}
\label{sec:pipeline}

The generated voices sounded pitched but blurred. To find out why, without
involving the model, we rendered 12-second excerpts of four held-out vocal
songs through each stage of the pipeline and compared the results with their
exact spectrograms (\cref{fig:errors,tab:reconstruction}).

\begin{figure}[htbp]
\centering
\resizebox{\linewidth}{!}{% Figure "errors", shared by the paper and (as SVG) the book.
\begin{tikzpicture}
\begin{axis}[
  xbar, bar width=9pt, width=72mm, height=46mm, scale only axis,
  xmin=0, xmax=1.1, xlabel={\sffamily\small mean error, log units per mel band},
  symbolic y coords={tokens to Griffin-Lim, tokenizer 4 levels, tokenizer 1 level, Griffin-Lim on exact mel},
  ytick=data, y dir=reverse, axis line style={draw=none}, tick style={draw=none},
  yticklabel style={font=\sffamily\small}, xticklabel style={font=\sffamily\scriptsize},
  xmajorgrids, grid style={white!85!black}, enlarge y limits=0.2,
  nodes near coords, nodes near coords style={font=\sffamily\scriptsize},
  every node near coord/.append style={xshift=1pt}]
\addplot[fill=accent!55, draw=none] coordinates {
  (0.194,Griffin-Lim on exact mel) (0.953,tokenizer 1 level)
  (0.704,tokenizer 4 levels) (0.710,tokens to Griffin-Lim)};
\end{axis}
\end{tikzpicture}
}
\caption{Mean spectral error of each stage on four held-out vocal excerpts,
for the first tokenizer (grey) and the final one (blue). The bottom pair is
the best that any model trained on these tokens could achieve.}
\label{fig:errors}
\end{figure}

\begin{widetable}
\caption{Spectral error of each stage on four held-out vocal excerpts, in log
units per mel band, for the first tokenizer and the final one.}
\label{tab:reconstruction}
\small
\begin{tabular}{lrrrrr}
\toprule
 & & \multicolumn{2}{c}{first tokenizer} & \multicolumn{2}{c}{final tokenizer} \\
\cmidrule(lr){3-4}\cmidrule(lr){5-6}
excerpt & vocoder & tokens & end to end & tokens & end to end \\
\midrule
guitar, brass, female voice & 0.230 & 0.807 & 0.821 & 0.425 & 0.480 \\
folk metal, male voice & 0.152 & 0.707 & 0.703 & 0.353 & 0.378 \\
electronic, vocal sample & 0.170 & 0.569 & 0.554 & 0.268 & 0.338 \\
bass, piano, wordless voice & 0.222 & 0.732 & 0.763 & 0.390 & 0.460 \\
\midrule
mean & 0.194 & 0.704 & 0.710 & 0.359 & 0.414 \\
\bottomrule
\end{tabular}
\end{widetable}

With the first tokenizer, the tokens alone caused 3.6 times the error of the
vocoder, and the end-to-end error of a perfect model, 0.71, was almost entirely
due to the tokenizer. An error of about 0.7 log units, or roughly 6\,dB per band,
is enough to blur the harmonic and formant detail on which an intelligible
voice depends. This measure does not capture the phase artefacts of
Griffin-Lim, which remain audible, so the tokenizer is the main spectral
limitation rather than the only audible one.

To examine the model on its own, we trained it on a single 12-second excerpt of
562 frames. Its training loss fell to 0.054 after 100 steps and to 0.002 after
200. We then generated 12 seconds from the excerpt's caption and compared the
generated codes with the excerpt's own (\cref{tab:memorisation}). Greedy
decoding reproduced all 562 frames, and top-$k$ sampling departed from the
excerpt only at isolated codes, without drifting away.

\begin{table}[htbp]
\centering
\caption{Generated codes compared with a memorised excerpt.}
\label{tab:memorisation}
\small
\begin{tabular}{lcc}
\toprule
decoding & agreement per codebook & identical prefix \\
 & (\%) & (frames) \\
\midrule
greedy & 100.0, 100.0, 100.0, 100.0 & 562 of 562 \\
top-$k$, seed 1 & 99.6, 100.0, 100.0, 99.6 & 26 \\
top-$k$, seed 2 & 100.0, 99.8, 99.3, 99.8 & 26 \\
\bottomrule
\end{tabular}
\end{table}

\begin{observation}
The model can reproduce audio exactly, and when it does, the result sounds
like the tokenizer's reconstruction. These two observations led us to
improve the tokens before training larger models.
\end{observation}

\subsection{A finer tokenizer}

Guided by this analysis, we compared tokenizers that differ in the number of
frames per step, the number of band groups and the number of levels
(\cref{tab:tokenizer-study}). Quantising one frame per step in 4 band groups
with 4 levels each lowered the mean error on the held-out songs from 0.776 to
0.401. A control without band groups suggests that most of this gain comes from
the larger number of codes per step, and a smaller part from the grouping. On the four vocal excerpts of \cref{fig:errors}, the error of the
tokens fell from 0.70 to 0.36, and the end-to-end error from 0.71 to 0.41. The
price is a vocabulary of 16 codebooks instead of 4.

\subsection{A larger model on the finer tokens}
\label{sec:large}

With the finer tokens we began training a larger model on the Thor: 8 blocks
of width 512, 42.9 million parameters, with the lyrics input enabled from the
start. It is trained on all 742 training songs of the combined corpora, with
the lyrics corpus weighted 4-fold and two takes recorded by the first author
weighted 10-fold, in batches of 8 sequences of 1\,408 positions. A step takes
2.0 seconds and about 11\,GB of memory. Training is still in progress; after
4\,500 of 20\,000 steps the held-out loss is 5.19 nats per code, averaged over
16 codebooks, and still falling. Because the vocabulary differs, this figure
is not comparable with those of \cref{tab:training}.

\subsection{Lyrics}

We extended both vocal models of \cref{tab:training} with a lyrics field and
trained them further on vocal songs, JamendoLyrics, and songs whose captions
mention rock or a related style. \Cref{tab:lyrics} gives the result of
\proc{Lyrics-Test} for a model trained on the Thor with the lyrics corpus
weighted 8-fold, over 48 batches of 4 windows.

\begin{table}[htbp]
\centering
\caption{Loss with the right lyrics, no lyrics and the wrong lyrics, and the
paired difference $\bar\Delta \pm s$.}
\label{tab:lyrics}
\small
\begin{tabular}{lrrrr}
\toprule
songs & right & none & wrong & $\bar\Delta \pm s$ \\
 & \multicolumn{4}{c}{(nats per code)} \\
\midrule
held out (6) & 4.3345 & 4.3409 & 4.3345 & $+0.0000 \pm 0.0004$ \\
training (73) & 4.0337 & 4.0395 & 4.0378 & $+0.0042 \pm 0.0006$ \\
\bottomrule
\end{tabular}
\end{table}

On the held-out songs, the right and the wrong lyrics give the same loss, so
this model does not yet use the words. On the training songs the difference is
statistically clear but small, which we take to indicate memorisation of those
songs rather than a learned relation between words and sound. We hope that the
paired test will be useful as a clear criterion for future attempts.

\section{Limitations}
\label{sec:limitations}

The tokenizer, although improved, still limits the fidelity of the output,
and Griffin-Lim adds phase artefacts that a learned vocoder would remove. The
data contain comparatively little singing (298 of the 706 songs mention a
voice) and very little of some styles (two songs are described as grunge).
Our models are much smaller than published systems \cite{copet2023musicgen},
and we have not yet carried out a formal listening test, so we make no claims
about perceived quality. On the implementation side, all arithmetic uses
32-bit floating point without tensor cores, and attention stores its full
matrix of scores \cite{dao2022flashattention}.

\section{Conclusion}
\label{sec:conclusion}

We have described StaleyNET, a complete text-to-music system written in Rust
and trained from scratch on embedded GPUs. Building every stage ourselves
allowed us to measure each one separately, and these measurements pointed to
the tokenizer, rather than the language model, as the main limitation on
sound quality. A revised tokenizer that quantises groups of frequency bands
roughly halved the reconstruction error, and a larger model is now being
trained on it. We hope that the worked examples, the stage-by-stage analysis
and the paired lyrics test will be of use to others who build similar systems
on modest hardware.

\FloatBarrier
